% =============================================================================
% 04-conduta.tex - PARTE 4, o que fazer agora. Frames 23 a 29.
% Farmacologia sempre em roxo.
% =============================================================================

% ----------------------------------------------------------------- frame 23 --
\begin{frame}{\ftit{\faKey}{Restri\c{c}\~ao de acesso a meios}}
\CRar{3}
\begin{center}
\begin{tikzpicture}
  \node[musialGreen, text width=9.4cm, font=\footnotesize, inner ysep=6pt] {%
    \itshape\guillemotleft Eu preciso conversar com voc\^es dois sobre deixar a
    casa mais segura nos pr\'oximos dias. N\~ao \'e desconfian\c{c}a. \'E que,
    quando aperta, a dist\^ancia at\'e a coisa faz diferen\c{c}a. O que a gente
    consegue tirar de casa hoje, e quem fica respons\'avel?\guillemotright};
\end{tikzpicture}
\end{center}

\CRvg
\begin{columns}[onlytextwidth, t]
\begin{column}{0.47\textwidth}
\begin{caixaproteger}{O que a conversa precisa produzir}
\begin{itemc}
\item o que sai de casa, nomeado
\item quem retira e para onde vai
\item at\'e quando fica assim
\item quem confere que foi feito
\end{itemc}
\end{caixaproteger}
\end{column}
\calha
\begin{column}{0.47\textwidth}
\begin{caixaagir}{O que invalida a medida}
\begin{itemc}
\item combinar s\'o com o paciente, e n\~ao com quem executa
\item deixar sem prazo e sem respons\'avel
\item esquecer a arma de fogo
\item n\~ao conferir depois
\end{itemc}
\end{caixaagir}
\end{column}
\end{columns}
\end{frame}

% ----------------------------------------------------------------- frame 24 --
\begin{frame}{\ftit{\faLifeRing}{Plano de seguran\c{c}a}}
\CRar{2.7}
\begin{columns}[onlytextwidth, T]
\begin{column}{0.50\textwidth}
\begin{center}
\begin{tikzpicture}[node distance=0.20cm]
  \node[crnode={campoambar}{3.55cm}, font=\scriptsize, inner ysep=3pt] (a)
    {\textbf{1.} sinais de alerta dele};
  \node[crnode={campoverde}{3.55cm}, font=\scriptsize, inner ysep=3pt,
        below=of a] (b) {\textbf{2.} estrat\'egias que ele j\'a usa};
  \node[crnode={campoverde}{3.55cm}, font=\scriptsize, inner ysep=3pt,
        below=of b] (c) {\textbf{3.} pessoas e ambientes que ajudam};
  \node[crnode={campoverde}{3.55cm}, font=\scriptsize, inner ysep=3pt,
        below=of c] (d) {\textbf{4.} quem pode ser acionado};
  \node[crnode={campoverde}{3.55cm}, font=\scriptsize, inner ysep=3pt,
        below=of d] (e) {\textbf{5.} servi\c{c}os e profissionais};
  \node[crnode={campobordo}{3.55cm}, font=\scriptsize, inner ysep=3pt,
        below=of e] (f) {\textbf{6.} redu\c{c}\~ao do acesso};
  \draw[crseta=campoverde] (a.south) -- (b.north);
  \draw[crseta=campoverde] (b.south) -- (c.north);
  \draw[crseta=campoverde] (c.south) -- (d.north);
  \draw[crseta=campoverde] (d.south) -- (e.north);
  \draw[crseta=campobordo] (e.south) -- (f.north);
\end{tikzpicture}
\end{center}
\end{column}
\calha
\begin{column}{0.44\textwidth}
\begin{caixaproteger}{N\~ao \'e contrato de n\~ao suic\'idio}
O contrato pede uma promessa e protege o servi\c{c}o. O plano descreve o que a
pessoa faz, na ordem, quando o pensamento aperta, e protege a pessoa. O
contrato n\~ao tem sustenta\c{c}\~ao e n\~ao substitui o plano.
\end{caixaproteger}
\CRvg
\begin{caixaproteger}{Escrito com ele}
Nas palavras dele, no celular dele, com os nomes e os n\'umeros preenchidos
ainda na consulta.
\end{caixaproteger}
\end{column}
\end{columns}

\fonte{Com contato de seguimento, 45\% menos comportamento suicida em seis meses.
Stanley et al., \emph{JAMA Psychiatry} 2018. \seloevidencia}
\end{frame}

% ----------------------------------------------------------------- frame 25 --
\begin{frame}{\ftit{\faStethoscope}{Tratar a crise e tratar a base}}
\CRar{1.5}
{\small O erro mais comum desta etapa \'e reduzir tudo ao antidepressivo e
deixar sem tratamento o que est\'a produzindo o risco hoje.\par}

\CRvg
\begin{columns}[onlytextwidth, t]
\begin{column}{0.47\textwidth}
\begin{caixamudanca}{O que tratar agora}
\begin{itemc}
\item intoxica\c{c}\~ao em curso
\item priva\c{c}\~ao de sono
\item dor n\~ao controlada
\item agita\c{c}\~ao
\item psicose
\item a crise psicossocial que precipitou
\end{itemc}
\end{caixamudanca}
\end{column}
\calha
\begin{column}{0.47\textwidth}
\begin{caixaproteger}{O que tratar em seguida}
\begin{itemc}
\item a condi\c{c}\~ao de base, com o esquema correto para ela
\item o que est\'a mal tratado ou sem ades\~ao
\item o que mudou de tratamento h\'a pouco
\item os fatores modific\'aveis que a fam\'ilia consegue sustentar
\end{itemc}
\end{caixaproteger}
\end{column}
\end{columns}

\mensagem{Cada fator modific\'avel tratado hoje \'e um modificador proximal a
menos na formula\c{c}\~ao de amanh\~a.}
\end{frame}

% ----------------------------------------------------------------- frame 26 --
\begin{frame}{\ftit{\faPills}{Fam\'ilia do f\'armaco e seguran\c{c}a da prescri\c{c}\~ao}}
\CRar{4.7}
\begin{caixafarmaco}{Seguran\c{c}a em superdosagem}
Perfis diferem quanto ao que acontece em ingest\~ao de toda a caixa. Isso entra
na escolha do f\'armaco quando h\'a suicidalidade ativa, e entra junto com
efic\'acia, comorbidade e ades\~ao, nunca sozinho.
\end{caixafarmaco}

\CRvg
\begin{caixafarmaco}{Quantidade dispensada}
Reduzir a quantidade entregue de uma vez \'e medida razo\'avel e de baixo custo
quando h\'a acesso e pouca supervis\~ao. \farmaco{N\~ao \'e crit\'erio isolado}:
n\~ao substitui a restri\c{c}\~ao de acesso a outros meios, n\~ao substitui a
supervis\~ao e n\~ao serve de justificativa para alta.
\end{caixafarmaco}

\CRvg
\begin{center}
\begin{tikzpicture}
  \node[crnode={campografite}{9.5cm}, font=\scriptsize] {%
    Prescrever para trinta dias a quem n\~ao tem supervis\~ao e tem
    idea\c{c}\~ao ativa \'e uma decis\~ao, ainda que ningu\'em a registre como
    tal.};
\end{tikzpicture}
\end{center}
\end{frame}

% ----------------------------------------------------------------- frame 27 --
\begin{frame}{\ftit{\faFlask}{Onde existe evid\^encia espec\'ifica}}
\CRar{4.8}
\begin{columns}[onlytextwidth, t]
\begin{column}{0.31\textwidth}
\begin{caixafarmaco}{L\'itio}
Em transtornos do humor, reduziu o n\'umero de suic\'idios em compara\c{c}\~ao
com placebo.\par
\farmaco{N\~ao \'e resgate.} O efeito \'e de manuten\c{c}\~ao, n\~ao de
plant\~ao.
\end{caixafarmaco}
\end{column}
\calha
\begin{column}{0.31\textwidth}
\begin{caixafarmaco}{Clozapina}
Em esquizofrenia e transtorno esquizoafetivo de alto risco, reduziu
comportamento suicida em compara\c{c}\~ao com olanzapina, em dois anos.
\end{caixafarmaco}
\end{column}
\calha
\begin{column}{0.31\textwidth}
\begin{caixafarmaco}{Escetamina}
Exemplo de redu\c{c}\~ao r\'apida de sintoma em contextos selecionados. No
Brasil, apenas a apresenta\c{c}\~ao nasal tem indica\c{c}\~ao aprovada; a
injet\'avel \'e uso off label.
\end{caixafarmaco}
\end{column}
\end{columns}

\CRvg
\mensagemagir{Melhora sintom\'atica r\'apida n\~ao equivale a preven\c{c}\~ao de
comportamento suicida.}

\fonte{Cipriani et al., \emph{BMJ} 2013 \seloevidencia \quad Meltzer et al.,
InterSePT, \emph{Arch Gen Psychiatry} 2003 \seloevidencia \quad Parecer CFM
20/2024 \selonorma}
\end{frame}

% ----------------------------------------------------------------- frame 28 --
\begin{frame}{\ftit{\faClock}{Janelas de vigil\^ancia}}
\CRar{4.8}
{\small Per\'iodos em que o risco sobe sem que o diagn\'ostico mude. Todos
previs\'iveis, todos registr\'aveis, todos ignorados com frequ\^encia.\par}

\CRvg
\begin{center}
\begin{tikzpicture}[node distance=0.26cm]
  \node[crnode={campoambar}{2.20cm}, font=\scriptsize, minimum height=1.00cm] (a)
    {alta recente};
  \node[crnode={campoambar}{2.20cm}, font=\scriptsize, minimum height=1.00cm,
        right=of a] (b) {in\'icio ou troca de tratamento};
  \node[crnode={campoambar}{2.20cm}, font=\scriptsize, minimum height=1.00cm,
        right=of b] (c) {agita\c{c}\~ao e ins\^onia};
  \node[crnode={campoambar}{2.20cm}, font=\scriptsize, minimum height=1.00cm,
        below=0.42cm of a, xshift=1.23cm] (d) {intoxica\c{c}\~ao};
  \node[crnode={campoambar}{2.20cm}, font=\scriptsize, minimum height=1.00cm,
        right=of d] (e) {mudan\c{c}a abrupta};
  \node[crnode={campoambar}{2.20cm}, font=\scriptsize, minimum height=1.00cm,
        right=of e] (f) {melhora aparente};
\end{tikzpicture}
\end{center}

\CRvg
\begin{caixamudanca}{A \'ultima \'e a mais trai\c{c}oeira}
A melhora aparente pode ser resolu\c{c}\~ao do quadro e pode ser a calma de quem
decidiu. As duas se parecem na sala. O que separa uma da outra \'e perguntar de
novo, e n\~ao presumir.
\end{caixamudanca}
\end{frame}

% ----------------------------------------------------------------- frame 29 --
\begin{frame}{\ftit{\faCalendarCheck}{Seguimento}}
\CRar{2.8}
\begin{center}
\begin{tikzpicture}[node distance=0.22cm]
  \node[crnode={campoverde}{1.70cm}, font=\scriptsize, inner xsep=4pt] (a) {data};
  \node[crnode={campoverde}{1.70cm}, font=\scriptsize, inner xsep=4pt, right=of a] (b) {local};
  \node[crnode={campoverde}{1.70cm}, font=\scriptsize, inner xsep=4pt, right=of b] (c) {respons\'avel};
  \node[crnode={campoverde}{1.70cm}, font=\scriptsize, inner xsep=4pt, right=of c] (d) {via de contato};
  \node[crnode={campobordo}{1.70cm}, font=\scriptsize, inner xsep=4pt, right=of d] (e) {conduta se piorar};
\end{tikzpicture}
\end{center}

\CRvg
\begin{caixaagir}{A express\~ao proibida}
\fala{Retorne se piorar} transfere ao paciente a tarefa de reconhecer a
piora e de decidir sozinho o que fazer com ela. Quem est\'a em crise \'e
exatamente quem menos consegue fazer isso.
\end{caixaagir}

\CRvg
\begin{caixaproteger}{O que entra no lugar}
\fala{Na quinta-feira, \`as 14h, na unidade de sa\'ude mental da Vila, com a
enfermeira Ana. Se antes disso os pensamentos voltarem como hoje, voc\^e liga
para este n\'umero, e sua irm\~a tamb\'em tem ele.}\par
\vspace{0.10cm}
Contato ativo do servi\c{c}o nas primeiras 48 a 72 horas, quando poss\'ivel.
\end{caixaproteger}
\end{frame}
